\section{Definición y tipos de sensores}

Un \textbf{sensor} es un dispositivo que detecta variaciones en una magnitud física y las convierte en señales útiles para un sistema de medición o control. En el contexto de IoT, los sensores permiten medir parámetros del entorno (luz, temperatura, humedad, etc.) y convertirlos en datos para la electrónica. Existen diversos tipos de sensores según la magnitud que miden y la naturaleza de su señal de salida. Por ejemplo:
\begin{itemize}
    \item \textbf{LDR (Light Dependent Resistor)}: sensor de luz ambiental. Es un fotoresistor pasivo cuya resistencia eléctrica varía inversamente con la intensidad luminosa incidente; a mayor iluminación, su resistencia disminuye, y viceversa \cite{Sensores-2}. Se utiliza típicamente en un divisor de voltaje para obtener una señal analógica proporcional a la luz.

    \item \textbf{Sensor de proximidad ultrasónico}: mide distancias mediante ultrasonido. Un módulo común (p. ej. HC-SR04) emite un pulso de sonido a alta frecuencia y espera el eco reflejado en un objeto; el tiempo transcurrido es proporcional a la distancia al obstáculo \cite{naylamp2025ultrasonico}. La distancia \(D\) se calcula usando la velocidad del sonido \(v\) (~340 m/s) y el tiempo de ida y vuelta \(t\): \(D = \frac{v \cdot t}{2}\). En la práctica, con \(t\) en microsegundos y \(D\) en cm, se usa la aproximación \(D,[cm] \approx t,[\mu s]/58\) \cite{Marquez2016Ultrasonico}.

    \item \textbf{Sensor de movimiento PIR (Passive Infrared)}: detecta presencia de personas u objetos calientes mediante cambios en la radiación infrarroja ambiental. Incorpora un elemento piroeléctrico que responde a variaciones en los niveles de infrarrojo emitido por el cuerpo humano u otras fuentes de calor en movimiento \cite{Sensores-1}. Entrega típicamente una salida digital (HIGH/LOW) indicando detección de movimiento.

    \item \textbf{Sensor de humedad}: higrómetro de suelo resistivo. Consta de dos electrodos insertados en la tierra; la conductividad entre ellos varía con la humedad del suelo. Un módulo conversor entrega una señal analógica inversamente proporcional a la humedad: en suelo seco la lectura es alta (~1023), y en suelo húmedo disminuye hacia 0 \cite{AbrahamG2020Humedad}.

    \item \textbf{Sensor de temperatura LM35}: entrega un voltaje linealmente proporcional a la temperatura (10 mV/°C). Con referencia de 5 V y resolución de 10 bits (Arduino), la temperatura \(T\) resulta de \(T(°C) = (5.0~\text{V} \cdot N / 1024) \times 100\), donde \(N\) es el valor leído (0–1023) \cite{Actuadores-2}.
\end{itemize}
Estos ejemplos ilustran sensores analógicos (LDR, LM35, higrómetro) y digitales (ultrasónico, PIR) utilizados en IoT. Cada sensor requiere un acondicionamiento adecuado de la señal para ser leído correctamente por un microcontrolador.

\section{Principios de funcionamiento de sensores analógicos y digitales}
La principal diferencia entre un sensor analógico y uno digital radica en la naturaleza de la señal de salida que producen \cite{Ruben2024Sensores}. Un sensor analógico proporciona una señal continua (típicamente un voltaje) que varía proporcionalmente con la magnitud física medida. Esta señal puede tomar cualquier valor dentro de un rango continuo. Por ejemplo, un termistor o un LDR pueden entregar cualquier voltaje entre 0 y 5 V dependiendo de la temperatura o luz incidente, incluyendo valores intermedios con decimales (ej. 2,72 V, 3,41 V, etc.) \cite{Sensores-1}. En cambio, un sensor digital produce una señal discreta, generalmente interpretada como estados lógicos. Muchos sensores digitales simplemente indican ON/OFF (1 o 0 lógico) según si la variable medida supera cierto umbral – por ejemplo, un PIR típicamente indica HIGH cuando detecta movimiento y LOW en reposo. Debido a esta naturaleza binaria, el resultado de un sensor digital no ofrece gradaciones intermedias de la magnitud medida, lo que puede limitar su resolución, mientras que los sensores analógicos, al tener salida continua, permiten una resolución más fina limitada solo por el convertidor analógico-digital (ADC) usado para leerlos \cite{Ruben2024Sensores}. Cabe mencionar que existen sensores digitales más complejos que comunican lecturas multibit (p. ej., sensores con interfaz $I^2C/SPI$ entregando valores cuantificados), pero internamente incluyen un ADC o lógica de conversión. En suma, sensores analógicos entregan señales continuas y requieren conversión ADC para ser interpretados numéricamente, mientras sensores digitales entregan directamente niveles lógicos o datos digitales ya procesados. Cada tipo tiene aplicaciones adecuadas: los analógicos se prefieren cuando se requiere granularidad en la medición, y los digitales cuando se necesitan umbrales sencillos o comunicaciones digitales robustas sin necesidad de calibrar un ADC.

\section{Actuadores y su interacción con microcontroladores}
Un actuador es, en términos generales, un dispositivo que recibe una energía de entrada (eléctrica, neumática, hidráulica, etc.) y la convierte en movimiento, fuerza u otra acción física útil \cite{Actuadores-2}. En sistemas IoT típicamente se emplean actuadores eléctricos para interactuar con el entorno, realizando acciones en base a las señales del sistema de control. Ejemplos de actuadores en IoT incluyen:
\begin{itemize}
    \item Motores de corriente continua (DC): convierten energía eléctrica en movimiento rotatorio. Son utilizados para mover robots, abrir/cerrar válvulas, bombear agua (una bomba de agua es esencialmente un motor DC acoplado a un mecanismo de bombeo), entre otros. Su velocidad de giro depende del voltaje aplicado o del ciclo de trabajo si se controlan por PWM (ver sección PWM) y el sentido de giro depende de la polaridad de alimentación.
    \item Servomotores y motores paso a paso: actuadores de posición controlada (no tratados en detalle aquí, pero comunes en robótica IoT).
    \item Solenoides y relés: producen movimiento lineal o acción de conmutación a partir de una bobina eléctrica (el relé se detalla más adelante como actuador de conmutación).
    \item Dispositivos de señalización: como LEDs (diodos emisores de luz) que convierten electricidad en luz, o buzzers que convierten señal eléctrica en sonido. Estos actuadores permiten retroalimentación al usuario (indicadores luminosos, alarmas sonoras).
\end{itemize}
Un aspecto crucial es cómo controla un microcontrolador a estos actuadores. Los pines de un microcontrolador (p.ej. Arduino) entregan señales lógicas de 5V o 3.3V con corriente limitada (típicamente ~20 mA máximo por pin para un Arduino UNO \cite{Arduino2024DigitalPins}). Esto es suficiente para activar directamente cargas pequeñas como un LED (colocando una resistencia en serie para limitar corriente) o un buzzer piezoeléctrico. Sin embargo, para actuadores que demandan mayor potencia – motores, bombas, solenoides – la corriente o voltaje requeridos exceden con creces lo que un pin puede proporcionar. Por ello, se emplean etapas de potencia entre el microcontrolador y el actuador. Estas etapas incluyen:
\begin{itemize}
    \item Transistores de conmutación (BJTs, MOSFETs) usados como interruptores electrónicos que permiten controlar corrientes grandes hacia el actuador mediante una señal de control de baja potencia proveniente del microcontrolador. Por ejemplo, un transistor MOSFET de canal N puede activar/desactivar un motor conectando este a una fuente externa de mayor corriente, recibiendo en su compuerta la señal de un pin digital.
    \item Drivers y puentes H integrados para motores (L293D, L298N, etc., discutidos en sección siguiente), que internamente contienen arreglos de transistores para controlar motores en ambos sentidos de giro y con señales PWM.
    \item Relés (electromecánicos o de estado sólido), que actúan como interruptores aislados galvánicamente: una bobina accionada por el microcontrolador conmuta contactos que pueden estar en un circuito de mayor voltaje/corriente. Esto permite encender/apagar cargas de alta potencia (ej. un motor alimentado con 12 V) sin conectar directamente dichas tensiones al microcontrolador \cite{Smoot2023Reles}.
\end{itemize}

\section{Modulación por Ancho de Pulso (PWM)}
La Modulación por Ancho de Pulso (PWM, Pulse Width Modulation) es una técnica electrónica fundamental para controlar la potencia entregada a un actuador mediante señales digitales. Consiste en generar una onda cuadrada (nivel alto y bajo) a una frecuencia determinada, y variar el ciclo de trabajo (duty cycle), es decir, el porcentaje del tiempo que la señal permanece en nivel alto durante cada período. Aunque la tensión instantánea solo toma valores altos o bajos, al promediarse en el tiempo la señal equivale a una tensión continua efectiva menor. Por ejemplo, un ciclo de trabajo del 50\% implica que el pin está en HIGH la mitad del tiempo y LOW la otra mitad; la tensión media en la carga será aproximadamente 50\% de la tensión de alimentación. Matemáticamente, si $V_{max}$ es el nivel alto (p. ej. 5 V) y el duty cycle es $\delta$ (en tanto por uno, o \%/100), la tensión media $V_{\text{avg}} \approx \delta \cdot V_{max}$. Igualmente, la potencia media entregada a la carga varía con ese factor $\delta$. 
\\\\
Dos parámetros caracterizan una señal PWM: la frecuencia de conmutación (Hz) y el duty cycle (\%). Es importante elegir una frecuencia suficientemente alta para que la inercia del sistema suavice la entrega de energía (por ejemplo, en motores o lámparas LED, una frecuencia de ~500 Hz o más evita parpadeos o vibraciones perceptibles). Por otro lado, el duty cycle es ajustado dinámicamente para regular la velocidad de motores, la intensidad lumínica de LEDs, etc. En aplicaciones de control de motores de corriente continua, un duty cycle bajo (p.ej. 10\%) entrega una energía promedio pequeña al motor, resultando en baja velocidad de giro, mientras que un duty cycle alto (90\%) suministra mayor energía promedio, permitiendo mayor velocidad. La modulación PWM es extremadamente eficiente, ya que los transistores de conmutación operan en corte o saturación (estado completamente off o on), minimizando disipación de potencia. Arduino dispone de salidas PWM dedicadas (pines ~3,5,6,9,10,11 en el UNO, más en el Mega) que usan temporizadores internos para generar estas señales. El programador puede establecer el duty cycle mediante la función analogWrite(pin, valor) en Arduino (donde valor de 0 a 255 representa 0\% a 100\% de ciclo útil). En síntesis, PWM permite simular un nivel analógico variable usando una señal digital, siendo esencial para controlar la velocidad de motores DC, la posición de servomotores (mediante pulsos PWM específicos), la intensidad de luces LED y en general cualquier actuador cuyo desempeño dependa de la potencia o voltaje aplicado. Su fórmula básica de control es: 
$Senal media = Duty Cycle \times Nivel Alto$. 
Aplicado a control de velocidad: 
$ \omega_{motor} = \alpha Duty Cycle \times \omega_{max} $, asumiendo comportamiento lineal en primera aproximación.

\section{Drivers de motor (L293D, L298N) – Topología Puente H}
Para controlar motores DC en ambos sentidos de giro mediante un microcontrolador, se utilizan circuitos denominados Puente H. Un puente H es una configuración de cuatro interruptores (transistores o MOSFETs) dispuestos en forma de “H” que permite invertir la polaridad aplicada a un motor: al cerrar dos transistores en diagonal se aplica voltaje en un sentido, y al conmutar al par opuesto se invierte la corriente por el motor, logrando el giro contrario. Integrar todos estos componentes discretos suele ser engorroso, por lo que existen drivers integrados como el L293D y L298N, que incorporan dos puentes H duales junto con protecciones, facilitando el control de dos motores DC (o un motor paso a paso bipolar) desde señales de lógica TTL. 
\\\\
L293D es un circuito integrado de puente H dual bastante difundido. Permite controlar dos motores DC en velocidad (mediante PWM) y dirección de giro de forma independiente \cite{Actuadores-1}. Soporta corrientes de hasta ~0.6 A continuas por canal (picos de 1.2 A) y voltajes de alimentación para motores desde ~4.5 V hasta 36 V \cite{Actuadores-1}. Internamente incluye diodos de protección (flyback) para manejar de forma segura cargas inductivas como motores \cite{Actuadores-1}. Cada canal del L293D tiene entradas de control (por lo general dos pines lógicos para decidir la dirección, y un pin enable que puede recibir PWM para regular la velocidad) \cite{Actuadores-1} \cite{Actuadores-2}. Una ventaja del L293D es que posee alimentación de lógica separada (5 V) y alimentación de potencia para motores, permitiendo usar fuentes externas adecuadas para la carga, aisladas de la electrónica de control (solo se debe compartir el GND común). En resumen, este chip actúa como intermediario: recibe señales de 5 V de Arduino en sus entradas, y conmuta una fuente de mayor corriente hacia el motor en la dirección indicada. Cada puente H entrega al motor una diferencia de potencial variable en sentido directo o inverso, posibilitando el giro en ambos sentidos \cite{Actuadores-1}. 
\\\\
L298N es otro driver de puente H dual, capaz de manejar corrientes mayores (hasta ~2 A por canal con disipador) y frecuentemente montado en módulos con bornes de conexión y diodos integrados. Opera de forma similar al L293D en términos lógicos: tiene entradas para cada medio puente (IN1, IN2, etc.) y entradas de enable (ENA, ENB) para activar cada canal y eventualmente alimentar PWM. Por ejemplo, en un módulo típico L298N, se conectan IN1/IN2 a pines digitales para controlar la dirección de un motor, y ENA a un pin PWM para controlar su velocidad\cite{Actuadores-2}. Al igual que el L293D, el L298N permite alimentar los motores con una fuente externa (hasta ~46 V según hoja de datos) separada de la lógica de 5 V. La topología interna del L298N consta de transistores Darlington en puente, lo que provoca una caída de voltaje relativamente alta en saturación; a pesar de ser un diseño veterano, sigue siendo muy usado por su facilidad de implementación. En la práctica, estos drivers simplifican enormemente la conexión: por ejemplo, para un motor se conecta sus terminales a las salidas OUT1/OUT2 del módulo, la fuente de poder del motor al pin VS (12 V, 5 V, etc. según requiera el motor) y el módulo se encargará de conmutar la energía al motor conforme a las señales de Arduino. Una pequeña tabla de verdad ilustra el control: si IN1=HIGH e IN2=LOW, el motor gira en un sentido; IN1=LOW, IN2=HIGH invierte el giro; si ambos IN1 e IN2 están en LOW (o HIGH) simultáneamente, el motor se frena (por cortocircuito dinámico o brake). La velocidad se ajusta modulando ENA con PWM (0-255 correspondie ndo a 0-100\% duty) \cite{Actuadores-2}. Estos drivers integrados, en esencia, implementan el puente H de forma compacta y protegen al microcontrolador de las altas corrientes y picos de voltaje generados por los motores.

\section{Relés electromecánicos}
Un relé es un actuador electromecánico que funciona como un interruptor controlado eléctricamente. Consiste en una bobina, un armadura móvil y contactos (normalmente abierto –NO–, normalmente cerrado –NC– y común –COM–). Cuando la bobina se energiza (con corriente proveniente del circuito de control), genera un campo magnético que atrae la armadura y cambia el estado de los contactos: el contacto NO se cierra con COM, y el NC se abre, o viceversa al desenergizar. La gran ventaja de un relé es que proporciona aislamiento galvánico entre la etapa de control y la etapa de potencia: la señal de baja potencia (ej. 5 V, 50 mA de Arduino vía un transistor) controla un circuito de alta potencia completamente separado (ej. 220 V AC o una fuente de 12 V DC para un motor)\cite{Smoot2023Reles}. Esto garantiza que transitorios o altos voltajes del lado de carga no afecten directamente al microcontrolador, aumentando la seguridad y protegiendo la electrónica de control. En IoT y sistemas embebidos, los relés se usan para accionar actuadores como motores de corriente alterna, luminarias, electrodomésticos o cualquier carga que exceda las capacidades electrónicas directas del MCU. 
\\\\
Al usar un relé con un microcontrolador, se debe considerar que la bobina del relé típicamente requiere más corriente de la que un pin digital puede suministrar (p. ej. 70-100 mA a 5 V, mientras un pin da ~20 mA). Por ello, se emplea un transistor de conmutación para excitar la bobina. Comúnmente se usan módulos de relé pre-diseñados que integran: un transistor de manejo, un diodo flyback (en paralelo con la bobina, para disipar el pico de voltaje cuando se desenergiza) y a veces un optoacoplador para aislamiento extra. La conexión típica involucra alimentar la bobina con 5 V (a Vcc del módulo), conectar su GND común al Arduino, y un pin digital activa la entrada IN del módulo (que a través del transistor energiza la bobina). Los contactos del relé entonces se usan en serie con la carga externa: por ejemplo, para controlar un motor DC con fuente separada mediante un relé, se conectaría el terminal COM del relé a la fuente positiva del motor, el terminal NO al terminal del motor, y el otro terminal del motor a la masa de la fuente. Cuando el Arduino activa el relé, los contactos COM-NO cierran y aplican la tensión de la fuente al motor, encendiéndolo \cite{Actuadores-2}. Al desactivar, el circuito se abre y el motor se apaga. Este método on/off es útil para controlar motores en un solo sentido sin modulación de velocidad, o para controlar cargas AC. Se sacrifica la rapidez (un relé conmutando mecánicamente tiene tiempos de respuesta del orden de milisegundos y limitada vida útil de ciclos), pero gana en aislamiento y capacidad de manejar altos voltajes/corrientes. En resumen, el relé actúa como un interruptor aislado accionado electrónicamente, imprescindible cuando se necesita separar galvanicamente un microcontrolador de una carga potente. Es un componente de interfaz muy común en IoT, por ejemplo en módulos de relé controlados por Arduino para domótica (encender luces de 220 V, motores de bombas de agua, etc.), garantizando que la etapa de control de baja potencia permanezca segura mientras gobierna dispositivos de alta potencia.

\section{Lectura analógica y digital en microcontroladores}
Los microcontroladores como Arduino UNO/Mega ofrecen pines de entrada digital y entrada analógica para conectar sensores. Las entradas digitales leen estados lógicos HIGH/LOW (1/0) basados en un umbral de voltaje (usualmente alrededor de 3 V para considerar HIGH con Vcc=5 V). Estas lecturas se manejan internamente con lógica booleana o variables enteras de valor 0/1. Por otro lado, las entradas analógicas están conectadas a un ADC (Analog-to-Digital Converter) interno, que convierte un voltaje analógico en un número digital discreto. En Arduino UNO/Mega, el ADC es de 10 bits de resolución, lo que produce un rango numérico de 0 a 1023 representando voltajes entre 0 V y la referencia seleccionada (usualmente 5 V por defecto)
\cite{Sensores-1}. Esto significa que la resolución es de ~4.9 mV por unidad (5 V/1023). La precisión de la medición viene dada por esta resolución junto con la tolerancia y ruido del sistema; en este caso, ~0.1\% del fondo de escala. Si se requiere mayor resolución (más bits) para lecturas muy precisas, podría usarse un microcontrolador con ADC de 12 bits, 16 bits o técnicas de oversampling; sin embargo 10 bits suele ser suficiente para la mayoría de aplicaciones de sensores comunes. 
\\\\
Al realizar una lectura analógica en Arduino (función analogRead(pin)), el valor retornado es un entero (int) entre 0 y 1023. Muchas veces este valor entero debe escalarse o convertirse a unidades físicas reales (°C, lux, etc.) mediante cálculos que involucran operaciones con decimales. Para ello, se emplean variables de coma flotante (float) en los cálculos. Por ejemplo, en la lectura del sensor LM35 mencionada, primero se obtiene el entero $N = \text{analogRead}(A0)$, luego se calcula la temperatura en ºC con la fórmula $T = (5.0 \times N \times 100.0)/1024.0$ \cite{Actuadores-2}. Aquí se usan constantes en punto flotante (5.0, 1024.0) y el resultado $T$ se guarda en una variable float para conservar los decimales (ya que la temperatura puede no ser entera). Si se usara solo enteros, la división truncaría cualquier parte fraccionaria, perdiendo precisión. En general, se usa tipo int para el valor bruto del ADC y tipo float para las magnitudes calibradas que requieran decimales. En cuanto a las lecturas digitales, normalmente se asignan a tipos lógicos o enteros pequeños ya que solo indican 0/1; no hay problema de precisión decimal allí. 
\\\\
Otra consideración es el rango de medida ajustable. Arduino permite cambiar la referencia de voltaje del ADC (por ejemplo a 1.1 V interna, 5 V default, o una referencia externa) para adaptarse al rango de salida del sensor y mejorar la resolución efectiva. Por ejemplo, si un sensor de temperatura LM35 nunca excederá 1 V de salida (100 ºC), usar la referencia interna de 1.1 V dará 1023 pasos en ~0–1.1 V, aumentando la sensibilidad a ~1 mV por unidad. También es posible implementar promedios o filtros en software para mejorar estabilidad en la lectura analógica, a costa de velocidad. 
\\\\
Finalmente, desde la perspectiva de programación, Arduino Mega2560 ofrece 16 canales de entrada analógica (A0–A15) y múltiples pines PWM, mientras Arduino Uno tiene 6 canales analógicos y 6 PWM. La elección entre int o float no solo afecta la precisión sino el uso de memoria y rendimiento: los cálculos en punto fijo (enteros) son más rápidos, pero cuando la aplicación requiere valores reales (ej. calcular humedad relativa, temperatura con decimales, tensión exacta de sensor), es imprescindible usar punto flotante. En resumen, la lectura digital es sencilla (alto/bajo) y se maneja con variables lógicas, la lectura analógica proporciona un número entero proporcional al voltaje de entrada, cuya correcta interpretación exige conocer el rango y resolución del ADC, y a menudo convierte ese entero a un valor físico usando operaciones en punto flotante para no perder detalle. Este manejo adecuado de tipos de datos garantiza que los sensores analógicos aprovechen al máximo la precisión disponible y que los cálculos en el microcontrolador reflejen fielmente las magnitudes medidas.